\subsection{Held-out confirmation, a network in the pool, and an expensive solver}\label{sec:confirm}
The criticisms that remained after the evaluation above were that only eight problems from four families were involved, that the tuned network may have been under-trained, that the default pool has no network, and that the solvers are so fast that the tool's cost was never weighed against a real saving. We therefore froze the tool, wrote a second protocol with criteria \emph{and predictions} (addendum v3.2; the freeze commit is recorded in Supplement, Sec.~\ref{sec:S-prereg}), and only then fitted anything to eight further \texttt{scikit-fem} problems (Supplement Table~\ref{tab:S-svproblems}; four new families: Helmholtz, nonlinear reaction--diffusion solved by Newton's method, elasticity on a disc, advection in a cellular flow). Unit tests of the new solvers exposed a boundary-region bug in one of them before any model was fitted (deviation V14). The pool, grids and criteria are those of the first evaluation. The problems were chosen for physical variety and two contain deliberately hard regimes; analyses added after seeing the results, in response to a simulated review, are labelled post hoc.
\paragraph{Held-out result} \Cref{tab:svcriteria} repeats the criteria of \cref{tab:hyp} on these eight problems; \SvMetCount{} of the \SvEvalCount{} evaluable verdicts were met (the basis flag never fired, so H3 could not be evaluated). Selection quality held up (H1a: median regret \SvSoneMedian, \SvSoneShare{} of cells within a factor two; the worst cell, \SvSthreeMaxWhere, has a cell regret of \SvSoneMax{} and the worst single job \SvWorstjob; \SvWorstjobOver{} of \SvWorstjobN{} jobs have a regret above two), and the tool was not worse than always using one Gaussian process (H1b, \SvStwoMedian). Four results went against the tool or against our predictions (\cref{tab:svpred}). (i) The forecast of H2 did not meet its criterion: \SvBoneK{} of \SvBoneN{} forecasts were within a factor two (\SvBoneShare; \SvBoneUnder{} were too optimistic by more than a factor two, median predicted/actual \SvBoneBias; from 64 runs the share was \SvBoneShareSixtyfour). The case-level interval \SvBoneLo{} to \SvBoneHi{} treats the 48 cases as independent; the problem-level interval, \SvPlBoneLo{} to \SvPlBoneHi, contains the threshold of 70\%, so a true rate near the threshold cannot be excluded with eight problems. Over all sixteen problems the share was \SvpBoneShare{} of \SvpBoneN{} (problem-level interval \SvpPlBoneLo{} to \SvpPlBoneHi); we describe the forecast as a rough guide that tends to be optimistic, not as validated. The classification of whether a target is reached was right in \SvBtwoOneAcc{} (1\%) and \SvBtwoTenthAcc{} (0.1\%) of cases, against \SvBtwoOneMajor{} and \SvBtwoTenthMajor{} for always predicting the majority class; a trivial rule that predicts ``reached'' whenever the cross-validated error at 256 runs is already below the target scores \SvCfOne{} and \SvCfTenth{} (post hoc), so the forecast adds little at 1\% and little at 0.1\%. (ii) The tool again did not match the benchmark's own selection (H1c: median ratio \SvSthreeMedian, better in \SvSthreeBetter{} cells and worse in \SvSthreeWorse; \SvSthreeWorseHelm{} of the worse cells are Helmholtz problems, where the benchmark's neural-network methods were selected and beat the tool; without the two Helmholtz problems the median ratio would be \SvNohSthree, better in \SvNohBetter{} and worse in \SvNohWorse{} of \SvNohN{} cells, post hoc). The two systems also differ in effort (five methods with eight search trials each against seven members with four-fold cross-validation), so H1c compares different tuning budgets in both directions. (iii) The basis flag never fired (H3 not evaluable), although in \SvHThreeFullBetter{} of the 72 fits a full-field model beat every POD-based member by more than 10\%: on these problems the flag has no recall. (iv) The tool is again at its limit against the tuned network (H4 \SvSfourVerdict: the network beat it by more than 25\% in \SvSfourK{} of \SvSfourN{} cells, median error ratio tool/network \SvSfourMedian). On the other side the reduced-basis group was clearly better at the largest budget on \SvPtK{} of the \SvPtN{} problems (H5 \SvPtVerdict; one of them, \texttt{disc\_elast\_L2}, is only just clear, $\bar y=\SvPtDisc$ against the threshold of 0.10), where we had predicted that it would not be. Problem-level 95\% intervals for the medians are \SvPlSoneLo{} to \SvPlSoneHi{} (H1a), \SvPlStwoLo{} to \SvPlStwoHi{} (H1b) and \SvPlSthreeLo{} to \SvPlSthreeHi{} (H1c).
\begin{table}[t]\centering\caption{Criteria on the first eight problems, on the eight further problems (tool frozen before they were run) and, as numbers only, pooled over sixteen. Each cell gives the observed value and the verdict; ``better/worse'' counts cells that differ by more than 10\%.}\label{tab:svcriteria}\footnotesize\resizebox{\textwidth}{!}{\begin{tabular}{p{1.3cm}p{4.4cm}p{3.4cm}p{3.6cm}p{3.4cm}}
\toprule
ID & Criterion & Stage 6 (8 problems) & Stage 7 (8 new, held out) & Pooled, 16 (numbers only) \\
\midrule
H1a & median regret $\le1.5$; $\ge75\%$ of cells $\le2$ & 1.00; 100\%; met & 1.04; 96\%; met & 1.01; 98\% \\
H1b & median ratio vs always-GP $\le1$ & 0.85; met & 0.97; met & 0.93 \\
H1c & median ratio vs benchmark selection $\le1$; better/worse cells & 1.01; 11/9; not met & 1.08; 6/10; not met & 1.02; 17/19 \\
H2 & forecast within $\times2$ in $\ge70\%$ & 77\%; met & 58\%; not met & 68\% \\
H2 (reach) & classification accuracy $\ge75\%$ (1\% / 0.1\%) & 100\% / 88\%; met & 100\% / 88\%; met & 100\% / 88\% \\
H3 & basis flag helps (if $\ge5$ flagged) & 15 flagged; met & 0 flagged; not evaluable & 15 flagged \\
H4 & tuned network better by $>25\%$ in $\le4$ of 16 & 4 of 16; met & 2 of 16; met & 6 of 32 \\
H5 & reduced basis clear on $\ge5$ of 8 & 3 of 8; not met & 6 of 8; met & 9 of 16 \\
\bottomrule
\end{tabular}
}\end{table}
\begin{table}[t]\centering\caption{Predictions written down in addendum v3.2 before the runs, and what happened (further eight problems).}\label{tab:svpred}\footnotesize
\begin{tabularx}{\textwidth}{>{\raggedright\arraybackslash}p{4.6cm}X}\toprule Prediction & Outcome\\\midrule
Selection (H1a) and the comparison with always-GP (H1b) hold & Held (median regret \SvSoneMedian; ratio \SvStwoMedian)\\
H1c marginal and may fail & Not met (\SvSthreeMedian). The addendum added that the criterion was ``only just met'' in the first evaluation; that was wrong, it was not met there either\\
Forecast holds from 64 runs, may fail from 256 & Failed from both (\SvBoneShareSixtyfour{} and \SvBoneShareTwofiftysix)\\
Classification of reaching a target will not beat the majority baseline by much & It beat it (\SvBtwoOneAcc{} and \SvBtwoTenthAcc), but a carry-forward rule does nearly as well\\
Reduced basis clearly better on fewer than five of eight (H5) & Better on \SvPtK{} of \SvPtN\\
Basis flag will probably fire on the resonant problem & It never fired; no recall (\SvHThreeFullBetter{} of 72 fits missed)\\
The resonant problem is the likeliest source of a large regret & The largest regret was on the other Helmholtz problem (\SvSthreeMaxWhere)\\
The stronger network will win more cells than the standard one & Four against three of sixteen\\
Tool + network meets W1; W2 uncertain & Both met; W2 is informative only in the two cells where the network changed the result\\
Expensive case: small regret, forecast uncertain & Regret \SvERegretTwoFiveSix; forecast within a factor two in \SvEFcK{} of \SvEFcN{} seeds (too pessimistic, the opposite direction)\\
\bottomrule\end{tabularx}\end{table}
\paragraph{The Helmholtz problems are a weak test} Both Helmholtz problems are undamped, and the range of the wave number contains an eigenvalue of the square (Supplement Table~\ref{tab:S-svproblems}), so the response has a pole inside the parameter range. The best pool member's relative test error is between \SvHelmMin{} and \SvHelmMax{} in every one of these six cells, that is, \emph{no model is usable}, and the regrets there compare models that all fail. The tool has no warning for this regime and prints a selection all the same; adding one (a statement when the best cross-validated error exceeds a stated level) is an obvious next step that we did not make after seeing these results. The largest regrets, the failure of H1c and most of the benefit of a network in the pool come from these two problems, so \cref{tab:svcriteria} should be read with that in mind.
\begin{table}[t]\centering\caption{Test error (mean relative $L_2$, geometric mean over three seeds) at $N=256$ on the eight further problems. ``Tool + net'' is the tool with a tuned network added to its pool.}\label{tab:sverrors}\small\resizebox{\textwidth}{!}{\begin{tabular}{lrrrrrr}
\toprule
Problem & Tool & Oracle & Always GP & Benchmark sel. & Tuned net & Tool + net \\
\midrule
\texttt{advdiff\_cell\_L1} & $1.0\times10^{-4}$ & $1.0\times10^{-4}$ & $2.9\times10^{-4}$ & $2.9\times10^{-4}$ & $2.9\times10^{-3}$ & $1.0\times10^{-4}$ \\
\texttt{advdiff\_cell\_L2} & $1.6\times10^{-3}$ & $1.4\times10^{-3}$ & $1.8\times10^{-3}$ & $1.8\times10^{-3}$ & $5.9\times10^{-3}$ & $1.6\times10^{-3}$ \\
\texttt{disc\_elast\_L1} & $2.8\times10^{-3}$ & $2.6\times10^{-3}$ & $2.1\times10^{-3}$ & $2.1\times10^{-3}$ & $8.7\times10^{-3}$ & $2.8\times10^{-3}$ \\
\texttt{disc\_elast\_L2} & $7.5\times10^{-2}$ & $7.4\times10^{-2}$ & $6.8\times10^{-2}$ & $6.8\times10^{-2}$ & $8.8\times10^{-2}$ & $7.5\times10^{-2}$ \\
\texttt{helmholtz\_L1} & $6.4\times10^{-1}$ & $5.8\times10^{-1}$ & $1.4\times10^{0}$ & $2.3\times10^{-1}$ & $2.7\times10^{-1}$ & $3.6\times10^{-1}$ \\
\texttt{helmholtz\_L2} & $1.0\times10^{0}$ & $9.7\times10^{-1}$ & $1.4\times10^{0}$ & $9.2\times10^{-1}$ & $2.0\times10^{0}$ & $1.0\times10^{0}$ \\
\texttt{reacdiff\_fe\_L1} & $1.4\times10^{-4}$ & $1.4\times10^{-4}$ & $1.6\times10^{-4}$ & $1.6\times10^{-4}$ & $1.2\times10^{-3}$ & $1.4\times10^{-4}$ \\
\texttt{reacdiff\_fe\_L2} & $4.3\times10^{-3}$ & $4.1\times10^{-3}$ & $4.4\times10^{-3}$ & $4.5\times10^{-3}$ & $3.9\times10^{-3}$ & $3.6\times10^{-3}$ \\
\bottomrule
\end{tabular}
}\end{table}
\paragraph{A network in the pool} Adding the tuned network to the pool (cross-validated like every other member, $N=256$, \cref{tab:sverrors}) changed the choice in \SvWSelNet{} of \SvWJobs{} fits and the result clearly in \SvWChanged{} of the eight problems (\SvWChangedNames; for \texttt{helmholtz\_L1} from \SvWHelmTool{} to \SvWHelmBoth, where the network alone gave \SvWHelmNet). The enlarged pool met both of its criteria: median regret \SvWoneMedian{} against the best member of the enlarged pool (largest \SvWoneMax, \SvWoneShare{} of cells within a factor two), and in \SvWtwoK{} of \SvWtwoN{} cells within 10\% of the better of the tool alone and the network alone. Because the network was not chosen in the other cells, the enlarged pool equals the default pool there, so W2 is informative only in the \SvWChanged{} cells that changed. The price is time: a median of \SvWSec{} seconds per fit, against about a hundred seconds for the default pool at this budget. The network is therefore offered as an option (\texttt{--with-network}), not as part of the default pool.
\paragraph{A larger tuning budget for the network} With three times the trials, depths up to five and widths up to 512, the network improved on the standard one in \SvSfiveBetter{} of \SvSfiveN{} problems at $N=256$ (median error ratio \SvSfiveRatio; \SvSfiveBetterOneZero{} improved by more than 10\%), at a median of \SvSfiveSec{} minutes per fit. It beat the tool by more than 25\% in \SvSfiveStr{} of the \SvSfiveN{} problems (\SvSfiveWins), against \SvSfiveStd{} for the standard network (\SvSfiveStdwins). In the median problem the tool's error was about half that of the stronger network (median ratio tool/network \SvSfiveToolratio). A larger search thus made the network somewhat better and did not change the picture: the network wins mostly on the level-two problems with more inputs or stronger nonlinearity, the tool elsewhere. We did not try larger budgets than these, so a network tuned further could do better.
\paragraph{An expensive solver, measured on a proxy} The L-shaped plate of the first evaluation was solved on a finely refined mesh, which takes \SvETsim{} seconds per run on this machine. This is the \emph{same problem as \texttt{lshape\_L1} on a finer mesh} (same inputs and outputs), so it tests the cost accounting, not selection or forecasting on a new problem, and it is still far from the hours-long simulations that motivate surrogates. 512 training runs and 300 test runs were computed once. Selection stayed close to the oracle (mean regret \SvERegretSixFour{} at 64 runs and \SvERegretTwoFiveSix{} at 256), and the error fell from \SvEErrSixtyfour{} to \SvEErrTwofiftysixPct{} between 64 and 256 runs. At 256 runs it was slightly worse than the best fixed member, \SvEBestfixed{} (\SvEBestfixedErr; the oracle in all three seeds), but better than always using the \texttt{matern52} Gaussian process (\SvEErrGpTwoFiveSix) and than the tuned network (\SvEErrNetTwoFiveSix, \SvENetSec{} seconds). The forecast from 64 to 256 runs was within a factor two in only \SvEFcK{} of \SvEFcN{} seeds; here it was too pessimistic (ratios \SvEFcRatioMin{} to \SvEFcRatioMax), the opposite of the optimistic bias on the held-out problems. Supplement Table~\ref{tab:S-svexp} gives the cost: by construction, the number of predictions after which the surrogate is cheaper than the solver is the number of training runs plus the tool time expressed in solver runs. The tool time in the table includes an evaluation-only step (refitting every pool member to score it on test data); without it the tool took \SvEToolsecNorefitTwofiftysix{} seconds at 256 runs (post hoc), and the break-even would be \SvEBreakNorefitTwofiftysix{} predictions. Ten thousand predictions would cost \SvEDirectH{} hours of solver time against \SvESurrH{} hours with the surrogate. A break-even is not an equal-accuracy comparison, so the error is given next to each value. This is one problem and one cost model; it neither includes validation runs nor parallel evaluation of the solver, and it is no evidence about simulations that take hours.
